% !PRINTED 138
\weberpage{138}

\begin{center}
{\Large Vierter Abschnitt.}\\[0.8em]
{\Large\gesperrt{Symmetrische Functionen.}}\\[0.8em]
\rule{0.16\textwidth}{0.4pt}
\end{center}

\begin{center}
{\large \S.\ 41.}

\vspace{0.5em}
{\large\gesperrt{Begriff der symmetrischen Functionen.}}

{\large\gesperrt{Symmetrische Grundfunctionen.}}
\end{center}

Wir betrachten in diesem Abschnitt \gesperrt{ganze Functionen} beliebigen Grades von einer beliebigen Anzahl $n$ von Veränderlichen
\[
 \alpha_1,\alpha_2,\ldots,\alpha_n.
\]

Eine solche Function heisst \gesperrt{symmetrisch}, \gesperrt{wenn sie ungeändert bleibt, wenn die Variablen $\alpha_1,\alpha_2,\ldots,\alpha_n$ einer beliebigen Permutation unterworfen werden,} und solche symmetrische Functionen sind es, deren Eigenschaften und Bildungsgesetze wir jetzt genauer kennen lernen müssen.

Damit eine Function $\Phi(\alpha_1,\alpha_2,\ldots,\alpha_n)$ symmetrisch sei, ist es genügend, dass sie sich bei der Vertauschung von je zweien der Argumente $\alpha_1,\alpha_2,\ldots,\alpha_n$ nicht ändere, weil nach \S.\,18 alle Permutationen durch eine Reihe von Transpositionen gebildet werden können.

Die Function $\Phi$ wird im Allgemeinen nicht homogen sein, sondern Glieder verschiedener Dimension enthalten; wenn man aber alle Glieder gleicher Dimension zusammenfasst, so lässt sich jedes $\Phi$ durch eine Summe \gesperrt{homogener Functionen} verschiedener Grade darstellen, und wenn $\Phi$ symmetrisch sein soll, so muss jeder homogene Bestandtheil eines bestimmten Grades für sich symmetrisch sein, da durch die Permutationen der Variablen die Dimensionen der Glieder nicht geändert werden.
